\documentclass[10pt,a4paper]{amsart}
\usepackage{amsmath,amssymb,amsthm,graphicx}
% For commutative diagrams we use Xy-pic
\usepackage[all]{xy}

% We use 2cell for 2-commutative diagrams.
\xyoption{2cell}
% Theorem environments.
%
\theoremstyle{plain}
\newtheorem{theorem}[subsection]{Theorem}
\newtheorem{proposition}[subsection]{Proposition}
\newtheorem{lemma}[subsection]{Lemma}

\theoremstyle{definition}
\newtheorem{definition}[subsection]{Definition}
\newtheorem{example}[subsection]{Example}
\newtheorem{exercise}[subsection]{Exercise}
\newtheorem{situation}[subsection]{Situation}

\theoremstyle{remark}
\newtheorem{remark}[subsection]{Remark}
\newtheorem{remarks}[subsection]{Remarks}

\numberwithin{equation}{subsection}

% Macros
%
\def\lim{\mathop{\mathrm{lim}}\nolimits}
\def\colim{\mathop{\mathrm{colim}}\nolimits}
\def\Spec{\mathop{\mathrm{Spec}}}
\def\Hom{\mathop{\mathrm{Hom}}\nolimits}
\def\Ext{\mathop{\mathrm{Ext}}\nolimits}
\def\SheafHom{\mathop{\mathcal{H}\!\mathit{om}}\nolimits}
\def\SheafExt{\mathop{\mathcal{E}\!\mathit{xt}}\nolimits}
\def\Sch{\mathit{Sch}}
\def\Mor{\mathop{\mathrm{Mor}}\nolimits}
\def\Ob{\mathop{\mathrm{Ob}}\nolimits}
\def\Sh{\mathop{\mathit{Sh}}\nolimits}
\def\NL{\mathop{N\!L}\nolimits}
\def\CH{\mathop{\mathrm{CH}}\nolimits}
\def\proetale{{pro\text{-}\acute{e}tale}}
\def\etale{{\acute{e}tale}}
\def\QCoh{\mathit{QCoh}}
\def\Ker{\mathop{\mathrm{Ker}}}
\def\Im{\mathop{\mathrm{Im}}}
\def\Coker{\mathop{\mathrm{Coker}}}
\def\Coim{\mathop{\mathrm{Coim}}}

% Boxtimes
%
\DeclareMathSymbol{\boxtimes}{\mathbin}{AMSa}{"02}

%
% Macros for moduli stacks/spaces
%
\def\QCohstack{\mathcal{QC}\!\mathit{oh}}
\def\Cohstack{\mathcal{C}\!\mathit{oh}}
\def\Spacesstack{\mathcal{S}\!\mathit{paces}}
\def\Quotfunctor{\mathrm{Quot}}
\def\Hilbfunctor{\mathrm{Hilb}}
\def\Curvesstack{\mathcal{C}\!\mathit{urves}}
\def\Polarizedstack{\mathcal{P}\!\mathit{olarized}}
\def\Complexesstack{\mathcal{C}\!\mathit{omplexes}}
% \Pic is the operator that assigns to X its picard group, usage \Pic(X)
% \Picardstack_{X/B} denotes the Picard stack of X over B
% \Picardfunctor_{X/B} denotes the Picard functor of X over B
\def\Pic{\mathop{\mathrm{Pic}}\nolimits}
\def\Picardstack{\mathcal{P}\!\mathit{ic}}
\def\Picardfunctor{\mathrm{Pic}}
\def\Deformationcategory{\mathcal{D}\!\mathit{ef}}
\usepackage{cleveref}
\crefname{originalremark}{original remark}{original remarks}
\Crefname{originalremark}{Original remark}{Original remarks}
\makeatletter
\newlabel{remark-more-elementary}{{02PK}{source}{}{sourcepointer-1}{}}
\newlabel{remark-more-elementary@cref}{{[originalremark][0][]02PK}{[1][0][]source}{}{}{}}
\makeatother
\title{The symbol of a discrete valuation ring equals the usual tame symbol}
\author{The Stacks Project Authors}
\date{stacks-project:chow.tex:local-snapshot-20261012; GNU FDL 1.2 or later}
\begin{document}
\maketitle
\noindent\textit{Editorial note.} All delivered mathematical text is the author's own, copied from the pinned Stacks Project chapter file chow.tex; no statement or proof step was rewritten, reordered or removed.\par\medskip
\noindent\textit{Editorial note.} Preamble provenance: the pinned chapter file opens with an input of the shared preamble and carries no preamble text of its own, so the project's single Stacks-wide preamble staged at sources/stacks-chow/preamble.tex is excerpted verbatim as two spans, lines 14-18 (Xy-pic setup) and lines 154-217 (the native theorem declarations and mathematical macros). Nothing was added, removed or redefined; the wrapper class and the amsmath/amssymb/amsthm/graphicx base packages are supplied by the assembler.\par\medskip
\noindent\textit{Editorial note.} The counted statement and proof contain no cross reference, no citation and no included file, so the delivered body is self contained.\par\medskip
\noindent\textit{Editorial note.} Canonical location of the counted proposition: Stacks tag 02Q6 (sources/stacks-chow/pinned-original-source.tex line 17624).\par\medskip
\noindent\textit{Allowed standard prerequisites (editorial).}\begin{itemize}
\item The symbol d\_A of a discrete valuation ring and its characterisation by one symbol comparison. Same chapter, chow.tex, the definition of d\_A in the tame symbol section and Remark labelled remark-more-elementary (source line 16632), invoked at proof line 17643. Supplies that d\_A(x, y) is the unique scalar making the displayed comparison of the two symbols hold; the choice of admissible sequences and the reduction to units are inside the counted span.
\item Multiplicativity of the symbol and the admissible relations for symbols. Same chapter, chow.tex, the multiplicativity statement used at proof line 17637 and the admissible relations used at proof lines 17651-17653. Supply the reduction to the case x, y in A and the final evaluation of the two symbols modulo the maximal ideal.
\end{itemize}\medskip
\ifdefined\hypertarget\hypertarget{sourcepointer-1}{}\fi
\noindent\textit{Source reference.} Original remark 02PK: sources/stacks-chow/pinned-original-source.tex, source lines 16632--16632. Same-chapter pointer into the author environment remark at source line 16632; invoked at proof line 17643 of the counted proof. It is retained only as a navigation pointer; every load-bearing step is inside the delivered span.\par\medskip
\begin{lemma}
\label{lemma-symbol-is-usual-tame-symbol}
Let $A$ be a discrete valuation ring with fraction field $K$.
For nonzero $x, y \in K$ we have
$$
d_A(x, y)
=
(-1)^{\text{ord}_A(x)\text{ord}_A(y)}
\frac{x^{\text{ord}_A(y)}}{y^{\text{ord}_A(x)}} \bmod \mathfrak m_A,
$$
in other words the symbol is equal to the usual tame symbol.
\end{lemma}

\begin{proof}
By multiplicativity it suffices to prove this when $x, y \in A$.
Let $t \in A$ be a uniformizer.
Write $x = t^bu$ and $y = t^bv$ for some $a, b \geq 0$
and $u, v \in A^*$. Set $l = a + b$. Then
$t^{l - 1}, \ldots, t^b$ is an admissible sequence in
$(x)/(xy)$ and $t^{l - 1}, \ldots, t^a$ is an admissible
sequence in $(y)/(xy)$. Hence by Remark \ref{remark-more-elementary}
we see that $d_A(x, y)$ is characterized by the equation
$$
[t^{l - 1}, \ldots, t^b, v^{-1}t^{b - 1}, \ldots, v^{-1}]
=
(-1)^{ab} d_A(x, y)
[t^{l - 1}, \ldots, t^a, u^{-1}t^{a - 1}, \ldots, u^{-1}].
$$
Hence by the admissible relations for the
symbols $[x_1, \ldots, x_l]$ we see that
$$
d_A(x, y) = (-1)^{ab} u^a/v^b \bmod \mathfrak m_A
$$
as desired.
\end{proof}

\end{document}
